% Einheit 1 von 4 · Wurzelziehen und Lösbarkeit (Katalog Einheit 1)
\hypertarget{e1}{}\einheitenkopf{Einheit 1 von 4 · Wurzelziehen und Lösbarkeit}
\zweigzeile{Hier lernst du, Gleichungen mit $x^2$ durch Wurzelziehen zu lösen und zu sagen, wie viele Lösungen sie haben · neu in diesem Jahr, je nach Buch erst Klasse 10 (am Gymnasium je nach Buch schon seit Klasse 8) · P10 · baut auf: Quadrieren, Quadratwurzeln, lineare Gleichungen, Einsetzen, Scheitelpunktform}
\renewcommand{\theaufgabe}{\arabic{aufgabe}}\setcounter{aufgabe}{0}

\begin{aufgabe}{Ich erkenne, ob eine Gleichung quadratisch oder linear ist. Kreuze an, ob $x^2$ oder eine quadrierte Klammer mit $x$ vorkommt. Du musst nichts lösen.}
\begin{geruest}
\gz{a}{$x^2 + 3 = 12$}{\kreuz{quadratisch}\kreuz{linear}}
\gz{b}{$4x - 1 = 7$}{\kreuz{quadratisch}\kreuz{linear}}
\gz{c}{$(x - 2)^2 = 5$}{\kreuz{quadratisch}\kreuz{linear}}
\gz{d}{$3 \cdot (x + 1) = 9$}{\kreuz{quadratisch}\kreuz{linear}}
\gz{e}{$2x = x^2$}{\kreuz{quadratisch}\kreuz{linear}}
\end{geruest}
\end{aufgabe}

\begin{aufgabe}{Ich erkenne, ob die rechte Seite positiv, null oder negativ ist. Sieh nur auf die rechte Seite und kreuze an. Du musst nichts rechnen.}
\begin{geruest}
\gz{a}{$x^2 = 17$}{\kreuz{positiv}\kreuz{null}\kreuz{negativ}}
\gz{b}{$x^2 = -8$}{\kreuz{positiv}\kreuz{null}\kreuz{negativ}}
\gz{c}{$5x^2 = 0$}{\kreuz{positiv}\kreuz{null}\kreuz{negativ}}
\gz{d}{$x^2 = 2{,}25$}{\kreuz{positiv}\kreuz{null}\kreuz{negativ}}
\gz{e}{$x^2 = -0{,}5$}{\kreuz{positiv}\kreuz{null}\kreuz{negativ}}
\end{geruest}
\end{aufgabe}

\begin{aufgabe}{Ich kann eine Gleichung der Form $x^2 = \text{Zahl}$ durch Wurzelziehen lösen. Gib beide Lösungen an.}
\beispiel{x^2 &= 25 &&\mid \text{Wurzel ziehen, beide Vorzeichen} \\ x_1 &= 5, \quad x_2 = -5}
\begin{gleichungsraster}[1]
\gl{x^2 = 4} & \gl{x^2 = 81} \\ \rule{0pt}{2mm} & \\
\gl{x^2 = 1} & \gl{x^2 = 100} \\ \rule{0pt}{2mm} & \\
\end{gleichungsraster}
\begin{teile}
\teil Ein quadratischer Garten hat den Flächeninhalt $64\,\text{m}^2$. Wie lang ist eine Seite?\\ Seitenlänge: \leerfeld[m]
\end{teile}
Löse. Geht die Wurzel nicht auf, lass sie stehen und runde auf zwei Stellen nach dem Komma.
\begin{gleichungsraster}[2]
\gl{x^2 = 11} & \gl{x^2 = -9} \\ \rule{0pt}{2mm} & \\
\end{gleichungsraster}
\end{aufgabe}

\begin{aufgabe}{Ich kann $x^2$ erst freistellen und dann die Wurzel ziehen. Löse die Gleichung.}
\beispiel{2x^2 - 32 &= 0 &&\mid +32 \\ 2x^2 &= 32 &&\mid :2 \\ x^2 &= 16 \\ x_1 &= 4, \quad x_2 = -4}
\begin{gleichungsraster}[3]
\gl{3x^2 = 75} & \gl{x^2 - 12 = 52} \\ \rule{0pt}{2mm} & \\
\gl{3x^2 - 7 = 5} & \gl{19 - x^2 = 3} \\ \rule{0pt}{2mm} & \\
\end{gleichungsraster}
\begin{teile}
\teil Eine zylinderförmige Dose soll $1000\,\text{cm}^3$ fassen und $12\,\text{cm}$ hoch sein. Stelle $V = \pi \cdot r^2 \cdot h$ nach $r^2$ um und berechne den Radius $r$. Runde auf zwei Stellen nach dem Komma. (P10 2022) \quad \feld[cm]{r}
\end{teile}
Stelle $x^2$ frei. Achte darauf, welche Zahl dann rechts steht.
\begin{gleichungsraster}[3]
\gl{2x^2 + 9 = 1} & \gl{3x^2 + 5 = 5} \\ \rule{0pt}{2mm} & \\
\end{gleichungsraster}
\end{aufgabe}

\begin{aufgabe}{Ich kann eine Gleichung mit quadrierter Klammer rückwärts lösen. Ziehe die Wurzel mit beiden Vorzeichen und löse dann die beiden linearen Gleichungen.}
\beispiel{(x - 3)^2 &= 25 \\ x - 3 &= 5 \quad\text{oder}\quad x - 3 = -5 &&\mid +3 \\ x_1 &= 8, \quad x_2 = -2}
\begin{gleichungsraster}[2]
\gl{(x - 1)^2 = 9} & \gl{(x - 2)^2 = 16} \\ \rule{0pt}{2mm} & \\
\gl{(x - 5)^2 = 1} & \gl{(x + 2)^2 = 9} \\ \rule{0pt}{2mm} & \\
\gl{(x - 4)^2 = 0} & \\ \rule{0pt}{2mm} & \\
\end{gleichungsraster}
\end{aufgabe}

\begin{aufgabe}{Ich kann durch Einsetzen prüfen, ob eine Zahl eine Lösung ist. Setze ein, rechne beide Seiten aus und kreuze an.}
\beispiel{x = -2 \text{ in } x^2 + 5 = 9{:}\quad (-2)^2 + 5 &= 9 \\ 4 + 5 &= 9 \\ 9 &= 9 && \text{(wA)}}
\begin{geruest}
\gz{a}{$x = 1$ \ in \ $x^2 + 3 = 4$}{\kreuz{(wA)}\kreuz{(fA)}}
\gz{b}{$x = 3$ \ in \ $x^2 - 1 = 10$}{\kreuz{(wA)}\kreuz{(fA)}}
\gz{c}{$x = 4$ \ in \ $2x^2 = 32$}{\kreuz{(wA)}\kreuz{(fA)}}
\gz{d}{$x = 5$ \ in \ $30 - x^2 = 5$}{\kreuz{(wA)}\kreuz{(fA)}}
\gz{e}{$x = -3$ \ in \ $2x^2 = 16$}{\kreuz{(wA)}\kreuz{(fA)}}
\gz{f}{$x = -1$ \ in \ $(x - 2)^2 = 1$}{\kreuz{(wA)}\kreuz{(fA)}}
\gz{g}{$x = -7$ \ in \ $(x + 3)^2 = 16$}{\kreuz{(wA)}\kreuz{(fA)}}
\end{geruest}
\end{aufgabe}

\begin{aufgabe}{Ich kann am Graphen ablesen, wie viele Lösungen eine Gleichung hat. Die Parabel gehört zu $y = (x - 1)^2 - 2$.}

\begin{ksys}[xmin=-2,xmax=6,ymin=-4,ymax=4,ablesen]
\parabel{1}{1}{-2}{}
\gerade[4.2]{0}{2}{y=2}
\gerade[4.2]{0}{-1}{y=-1}
\gerade[4.2]{0}{-2}{y=-2}
\end{ksys}

Beispiel: Die Gerade $y = 2$ schneidet die Parabel an zwei Stellen, bei $x = -1$ und bei $x = 3$. Also hat $(x - 1)^2 - 2 = 2$ zwei Lösungen: $x_1 = -1$, $x_2 = 3$.\par
Gib an, wie viele Lösungen die Gleichung hat, und lies die Lösungen ab.
\begin{geruest}
\gz{a}{$(x - 1)^2 - 2 = -1$}{Anzahl: \leerfeld \quad Lösungen: \leerfeld}
\gz{b}{$(x - 1)^2 - 2 = -2$}{Anzahl: \leerfeld \quad Lösungen: \leerfeld}
\gz{c}{$(x - 1)^2 - 2 = -3$}{Anzahl: \leerfeld \quad Lösungen: \leerfeld}
\end{geruest}
Entscheide ohne Zeichnung mit dem Scheitel, wie viele Lösungen die Gleichung hat.
\begin{geruest}
\gz{d}{$(x + 2)^2 + 1 = 0$}{Anzahl: \leerfeld}
\gz{e}{$-(x - 1)^2 + 3 = -2$}{Anzahl: \leerfeld}
\gz{f}{$(x - 3)^2 - 5 = -5$}{Anzahl: \leerfeld}
\end{geruest}
\end{aufgabe}

\begin{aufgabe}{Ich kann Gleichungen mit keiner, einer oder zwei Lösungen angeben und Aussagen dazu beurteilen. Ergänze die rechte Seite.}
\begin{teile}
\teil Die Gleichung soll keine Lösung haben: \quad $x^2 =$ \leerfeld
\teil Die Gleichung soll genau eine Lösung haben: \quad $x^2 + 7 =$ \leerfeld
\teil Entscheide, ob die Aussagen wahr oder falsch sind. Gib außerdem eine quadratische Gleichung an, die keine Lösung hat. (P10 2021)
\end{teile}
\sachtabelle{lcc}{Aussage & wahr & falsch}{$(x - 6)^2 = 0$ hat genau eine Lösung. & $\square$ & $\square$ \\ $2$ und $8$ sind Lösungen von $(x + 5)^2 = 9$. & $\square$ & $\square$}
\par\medskip Gleichung ohne Lösung: \leerfeld \leerfeld
\end{aufgabe}

\begin{aufgabe}{Ich finde den Fehler beim Wurzelziehen.}
\begin{teile}
\teil Ole löst $3x^2 = 48$ so:
\rechnung{3x^2 &= 48 &&\mid :3 \\ x^2 &= 16 \\ x &= 4}
Finde den Fehler, benenne ihn und korrigiere die Rechnung.
\end{teile}
Löse selbst.
\begin{gleichungsraster}[3]
\gl{2x^2 = 162} & \\ \rule{0pt}{2mm} & \\
\end{gleichungsraster}
\begin{teile}
\teil Lea löst $(x + 6)^2 = 4$ so:
\rechnung{(x + 6)^2 &= 4 \\ x + 6 &= 2 \quad\text{oder}\quad x + 6 = -2 \\ x_1 &= 8, \quad x_2 = 4}
Finde den Fehler, benenne ihn und korrigiere die Rechnung.
\end{teile}
Löse selbst.
\begin{gleichungsraster}[2]
\gl{(x + 5)^2 = 4} & \\ \rule{0pt}{2mm} & \\
\end{gleichungsraster}
\end{aufgabe}

\begin{aufgabe}{Ich kann begründen, wie viele Lösungen eine Gleichung mit $x^2$ hat.}
\begin{teile}
\teil Begründe, warum $x^2 = 30$ zwei Lösungen hat.
\teil Begründe ohne Rechnung, warum $x^2 = -30$ keine Lösung hat.
\teil Nina sagt: „$(x - 2)^2 = 0$ hat keine Lösung, weil man aus null keine Wurzel ziehen kann.“ Hat sie recht? Begründe.
\end{teile}
\end{aufgabe}
